%------------------------------------------------------------------------------%

\usepackage{apresentacao}
\usepackage{tabto}

\makeatletter
\graphicspath  {{./figs/}}
\def\input@path{{./figs/}}
\makeatother

%------------------------------------------------------------------------------%
\title   {Método Matemático}
\author  {\large Luis Alberto D'Afonseca}
\date    {Integrais e Séries}

%------------------------------------------------------------------------------%
\begin{document}

%------------------------------------------------------------------------------%
\begin{frame}
    \titlepage
\end{frame}

%------------------------------------------------------------------------------%
\section{O Método Matemático}

%------------------------------------------------------------------------------%
\begin{frame}{O Método (Secreto) da Matemática}
  \begin{enumerate}[<+->]
    \item Não basta encontrar a solução correta, \\ queremos a solução \emph{comprovadamente correta}
    \item Uma aproximação não é uma solução
    \item Não podemos usar o método empírico
    \item Toda afirmação precisa ser \emph{comprovadamente verdadeira}
    \begin{enumerate}
      \item Axiomas -- Assumimos que é verdade
      \item Proposições, Lemas, Teoremas, Corolários -- Demonstramos que é verdade
      \item Nem sempre o \textit{``obvio''} é verdadeiro
    \end{enumerate}
    \item Cada passagem em um desenvolvimento deve \emph{obrigatoriamente}
    estar ancorada em um resultado demonstrado (Por isso eles são numerados)
  \end{enumerate}
\end{frame}

%------------------------------------------------------------------------------%
\begin{frame}
  \tableofcontents
\end{frame}

%------------------------------------------------------------------------------%
\end{document}
%------------------------------------------------------------------------------%
